% ============================================================================
% TopK Sparse Autoencoders Across Three Model Architectures
%
% Scope note (2026-07-30 revision): this document previously described a
% three-way comparison of SAE training objectives (L1, TopK, Gated) across
% GPT-2 Small, Pythia-1.4B and Llama-3.2-3B, with results stated in the
% abstract and every results table left as a \placeholder. Only TopK SAEs
% were ever trained, on Llama-3.2-3B / Mistral-7B / Qwen2.5-3B, and no probe,
% human-rating or steering study was conducted. The paper has been rescoped to
% what the data supports. See \S\ref{sec:notrun} for what was planned and not
% run.
%
% Build: tectonic -X compile main.tex
% ============================================================================

\documentclass[11pt]{article}

\usepackage[margin=1.25in, top=1in, bottom=1.25in]{geometry}
\usepackage[utf8]{inputenc}
\usepackage[T1]{fontenc}
\usepackage{amsmath}
\usepackage{amssymb}
\usepackage{booktabs}
\usepackage{graphicx}
\usepackage{microtype}
\usepackage{xcolor}
\usepackage{url}
\usepackage[colorlinks=true, linkcolor=blue!70!black,
            citecolor=blue!70!black, urlcolor=blue!70!black]{hyperref}
\usepackage[numbers, sort&compress]{natbib}

\title{TopK Sparse Autoencoders Across Three Model Architectures:\\
Dictionary Collapse, Dense-Feature Degeneracy, and the\\
Limits of Activation-Pattern Feature Matching}

\author{J.~Melton \\
  Independent Researcher \\
  \texttt{jacob@paydirt.solutions}}

\date{July 2026 \\ \smallskip \small Preprint.}

\begin{document}
\maketitle

% ── Abstract ────────────────────────────────────────────────────────────────
\begin{abstract}
We train TopK sparse autoencoders (SAEs) at matched sparsity ($k = 128$) and matched
dictionary size ($16{,}384$) on mid-network residual stream activations from three
model families---Llama-3.2-3B, Mistral-7B-v0.3 and Qwen2.5-3B---and report three
findings, two of them negative.

First, TopK training without an auxiliary revival loss undergoes an abrupt
\emph{dictionary collapse}. In both runs long enough to exhibit it, dead features
roughly quadruple between step 2,000 and step 5,000 (Llama: 1,078 $\to$ 6,857; Qwen:
2,893 $\to$ 11,178) and never recover, because a feature below the top-$k$ threshold on
every batch receives exactly zero encoder gradient. Reconstruction quality is unharmed
and in fact continues improving throughout, so the collapse is invisible to the metric
practitioners usually monitor. Qwen ends at 57.4\% of its dictionary permanently dead.

Second, the surviving code is dominated by a small set of dense directions. With
$k/D = 0.78\%$, a feature should fire on under 1\% of tokens if activation were spread
evenly; instead 20--34 features per model fire on more than half of all tokens, and the
200 most frequent features absorb 39--57\% of total activation mass. Any procedure that
selects features by activation frequency---a common first step in feature labelling and
in feature-based classifiers---selects almost entirely from this degenerate head.

Third, cross-architecture feature matching by activation-pattern similarity finds
nothing. Under a chunk-permutation null, the best-of-16{,}384 match similarity between
\emph{unrelated} features averages 0.39--0.57 and has a 99th percentile of 0.83--0.94,
so the conventional fixed threshold of 0.8 sits at or below the noise floor. With
one-to-one reciprocal matching, a closed-triangle requirement and a null-calibrated
threshold, the number of features shared across all three models is zero; an earlier
uncalibrated analysis of the same checkpoints reported 3,753.

We report a comparison of training objectives as \emph{not conducted}: only TopK was
trained, and \S\ref{sec:notrun} states plainly which planned experiments do not exist.
\end{abstract}

% ── 1. Introduction ──────────────────────────────────────────────────────────
\section{Introduction}
\label{sec:intro}

Language models represent concepts in high-dimensional activation spaces that mix
thousands of distinct semantic and syntactic features within individual neurons---the
phenomenon \citet{elhage2022superposition} term \emph{superposition}. Sparse
autoencoders address this by learning an overcomplete dictionary of directions such
that any activation can be reconstructed as a sparse sum over dictionary elements.
\citet{bricken2023monosemanticity} applied the approach at scale and found the
resulting features highly interpretable.

The practical literature has converged on a small number of quality metrics---variance
explained, $L_0$, dead-feature count---and on reporting them at the end of training.
This paper is about what those metrics miss. We train TopK SAEs \citep{gao2024scaling}
under matched hyperparameters on three model families and instrument the training
process rather than only its endpoint. Two pathologies emerge that a final-checkpoint
report would not surface, and a third result---an attempt to match features across the
three models---fails in a way that we think is instructive.

\paragraph{Why these three models.} Llama-3.2-3B, Mistral-7B-v0.3 and Qwen2.5-3B come
from three organizations, have different hidden dimensions (3{,}072 / 4{,}096 /
2{,}048), different attention schemes, and no publicly shared training data. Holding
dictionary size fixed at 16{,}384 across them means the expansion ratio varies
(5.33$\times$ / 4.0$\times$ / 8.0$\times$), which turns out to matter: the model with
the largest expansion ratio collapses hardest.

\paragraph{Contributions.}
\begin{enumerate}
  \item \textbf{Dictionary collapse is abrupt, mid-training, and invisible to FVE.} We
    characterize a sharp transition between steps 2,000 and 5,000 in which dead
    features quadruple while reconstruction quality continues to improve
    (\S\ref{sec:collapse}).
  \item \textbf{Quantification of dense-feature degeneracy.} We measure how much of the
    sparse code's activation mass is carried by its most frequent features, and show
    the concentration is severe enough to invalidate frequency-based feature selection
    (\S\ref{sec:dense}).
  \item \textbf{A permutation null for cross-model feature matching, and a negative
    result.} We show that the standard protocol---fixed cosine threshold, nearest-neighbour
    matching---produces thousands of spurious matches, and that nothing survives proper
    calibration (\S\ref{sec:matching}).
  \item \textbf{A held-out generalization measurement} for the one SAE that has unseen
    data, and an explicit statement of why the other two do not (\S\ref{sec:heldout}).
\end{enumerate}

% ── 2. Related Work ──────────────────────────────────────────────────────────
\section{Related Work}
\label{sec:related}

\subsection{Superposition and Dictionary Learning}

\citet{elhage2022superposition} established the theoretical basis for SAE-based feature
extraction: a transformer with $n$ neurons can represent $O(n^2)$ nearly orthogonal
feature directions if those features are sparse. Dictionary learning recovers them by
sparse factorization: given activations $X \in \mathbb{R}^{n \times d}$, find
$D \in \mathbb{R}^{m \times d}$ (with $m \gg d$) and sparse $C \in \mathbb{R}^{n \times m}$
such that $X \approx CD$ and $\|c_i\|_0 \ll m$ for each row.

\subsection{SAE Architecture Variants}

\paragraph{Standard L1 SAE.}
\citet{cunningham2023sparse,bricken2023monosemanticity} train an encoder--decoder pair
to minimize
\begin{equation}
  \mathcal{L} = \|x - \hat{x}\|_2^2 + \lambda \|z\|_1 ,
  \label{eq:l1}
\end{equation}
where $\hat{x} = W_{\text{dec}} \operatorname{ReLU}(W_{\text{enc}} x + b_{\text{enc}})
+ b_{\text{dec}}$. The L1 penalty encourages sparsity through a continuous relaxation.

\paragraph{TopK SAE.}
\citet{gao2024scaling} replace the L1 penalty with a hard constraint:
\begin{equation}
  z^{\text{TopK}} = \operatorname{TopK}(\operatorname{ReLU}(W_{\text{enc}} x
    + b_{\text{enc}})) ,
  \label{eq:topk}
\end{equation}
guaranteeing exactly $k$ nonzero activations per input. Critically for this paper,
\citet{gao2024scaling} pair TopK with an \emph{auxiliary loss} that periodically
resurrects dead feature directions. Our implementation omits it, and
\S\ref{sec:collapse} is in effect a measurement of what that omission costs.

\paragraph{Gated and JumpReLU SAEs.}
\citet{rajamanoharan2024gated} separate feature detection from magnitude estimation via
a gating network; \citet{rajamanoharan2024jumprelu} use a discontinuous activation with
a learned threshold. We did not train either variant (\S\ref{sec:notrun}).

\subsection{Evaluation Methodologies}

Prior evaluation approaches include linear probing
\citep{alain2016understanding,belinkov2022probing,cunningham2023sparse}, human
interpretability ratings via the top-$k$ activating contexts protocol
\citep{bricken2023monosemanticity}, and activation steering fidelity
\citep{turner2023activation,zou2023representation}. This paper uses none of them; our
evaluation is confined to reconstruction quality, dictionary utilization, and
cross-model feature correspondence.

% ── 3. Methods ───────────────────────────────────────────────────────────────
\section{Methods}
\label{sec:methods}

\subsection{Models, Layers, and Activation Collection}

SAEs are trained on residual stream activations at the 50\% depth midpoint of each
model: \textbf{Llama-3.2-3B} (layer 14 of 28, $d = 3{,}072$), \textbf{Mistral-7B-v0.3}
(layer 16 of 32, $d = 4{,}096$), and \textbf{Qwen2.5-3B} (layer 18 of 36,
$d = 2{,}048$). Activations are collected from 500{,}000 tokens of WikiText-103
(\texttt{wikitext-103-raw-v1}, train split), streamed in the same order for every
model, and stored as float16.

Mistral activations were extracted from 4-bit quantized weights
(\texttt{mlx-community/Mistral-7B-v0.3-4bit}); the bf16 base model requires
authentication we did not have. This is a confound we cannot rule out and flag wherever
Mistral results appear.

\subsection{SAE Configuration}

All three SAEs are TopK with $k = 128$, dictionary size $D = 16{,}384$, peak learning
rate $10^{-4}$ with cosine annealing, batch size 2{,}048 tokens, and seed 42. Because
$D$ is held fixed while $d$ varies, the expansion ratio differs by model: 5.33$\times$
(Llama), 4.0$\times$ (Mistral), 8.0$\times$ (Qwen). The encoder is
$z = \operatorname{TopK}(\operatorname{ReLU}(W_{\text{enc}}(x - b_{\text{dec}})
+ b_{\text{enc}}))$ and the decoder $\hat{x} = W_{\text{dec}} z + b_{\text{dec}}$, with
decoder columns unit-normalized. The loss is plain MSE; there is no auxiliary
dead-feature loss and no re-initialization scheme.

Training runs differ in length, and not by design (Table~\ref{tab:runs}). Only the Qwen
run reached its 50{,}000-step target. This unevenness limits cross-model comparison
throughout, and we do not attribute differences between models to architecture where
training length is a plausible alternative explanation.

\begin{table}[t]
  \centering
  \caption{The three training runs. Only Qwen reached the step target. ``Epochs'' is
    the number of passes over the source token set implied by steps $\times$ batch.}
  \label{tab:runs}
  \smallskip
  \begin{tabular}{llrrrrr}
    \toprule
    Model & Layer & $d$ & Expansion & Source tokens & Steps & Epochs \\
    \midrule
    Llama-3.2-3B      & 14 & 3{,}072 & 5.33$\times$ & 500{,}000 & 13{,}000 & 53 \\
    Mistral-7B$^\ddagger$ & 16 & 4{,}096 & 4.0$\times$  &  50{,}000 &  1{,}000 & 41 \\
    Qwen2.5-3B        & 18 & 2{,}048 & 8.0$\times$  & 500{,}000 & 50{,}000 & 205 \\
    \bottomrule
  \end{tabular}

  \smallskip
  \noindent\footnotesize $\ddagger$~4-bit quantized source model.
\end{table}

\subsection{Reconstruction Metrics}

We report fraction of variance explained,
$\text{FVE} = 1 - \operatorname{Var}(x - \hat{x}) / \operatorname{Var}(x)$, and MSE.
MSE is activation-scale-dependent and is not comparable across models with different
hidden dimensions and activation norms; FVE is.

An important measurement detail: during training, FVE is computed on the current batch.
This is a noisy estimator, and reporting the final-step value---the natural thing to
do---samples once from a wide distribution. Across the last ten logged steps the batch
FVE spans 0.013 (Llama), 0.021 (Qwen) and 0.364 (Mistral). All FVE figures in
\S\ref{sec:results} are instead computed over the full 500{,}000-token corpus.

\subsection{Cross-Architecture Feature Matching}
\label{sec:matching-method}

Features are compared across models by \emph{functional} similarity: the correspondence
of their activation patterns over a shared corpus. For each model we encode 500{,}000
tokens of WikiText-103, partition them into $N = 1{,}000$ non-overlapping chunks of 500
tokens, and average post-TopK feature activations within each chunk. Column $j$ of the
resulting $N \times D$ matrix is a 1{,}000-dimensional \emph{fingerprint} of feature
$j$. Fingerprints are centred before normalization, so the similarity statistic is a
Pearson correlation across chunks rather than a cosine between non-negative vectors.

Three procedural choices distinguish this from the conventional protocol, and
\S\ref{sec:matching} shows each is necessary:

\begin{enumerate}
  \item \textbf{Support filter.} Features active in fewer than 5\% of chunks are
    excluded. A feature active in a handful of chunks correlates near 1 with any other
    feature active in the same chunks, by construction rather than by shared function.
  \item \textbf{Reciprocal matching.} A match requires features $i$ and $j$ to be
    \emph{mutual} nearest neighbours. This is one-to-one by construction. Taking the
    union of both directions' nearest neighbours---the common alternative---permits one
    feature to be the recorded partner of arbitrarily many, and permits match counts
    exceeding the dictionary size.
  \item \textbf{Null-calibrated threshold.} $\tau$ is set to the 99th percentile of a
    permutation null in which chunk order is shuffled in one model, destroying
    cross-model correspondence while preserving each feature's marginal activation
    profile (20 replicates).
\end{enumerate}

A feature triple is \emph{three-way universal} only if all three pairwise edges
independently satisfy these criteria---a closed triangle, not a chain through an
intermediate model.

% ── 4. Results ───────────────────────────────────────────────────────────────
\section{Results}
\label{sec:results}

\subsection{Dictionary Collapse}
\label{sec:collapse}

Table~\ref{tab:collapse} tracks dead features---those with no activation over a rolling
5{,}000-step window---through training.

\begin{table}[t]
  \centering
  \caption{Dead-feature count (of 16{,}384) over training. Both runs long enough to
    show it undergo a sharp transition between steps 2{,}000 and 5{,}000. FVE is the
    batch estimate at that step; note that it does not degrade across the transition.
    Data: \texttt{data/sae-runs/*/training.jsonl}.}
  \label{tab:collapse}
  \smallskip
  \begin{tabular}{rrrrrrr}
    \toprule
    & \multicolumn{2}{c}{Llama-3.2-3B} & \multicolumn{2}{c}{Qwen2.5-3B}
    & \multicolumn{2}{c}{Mistral-7B} \\
    \cmidrule(lr){2-3}\cmidrule(lr){4-5}\cmidrule(lr){6-7}
    Step & Dead & FVE & Dead & FVE & Dead & FVE \\
    \midrule
        1 &  1{,}381 & $-0.109$ &  3{,}220 & $-0.096$ & 748 & $-0.088$ \\
      200 &  1{,}245 &   0.624  &  3{,}091 &   0.743  & 468 &   0.829  \\
    1{,}000 &  1{,}145 &   0.955  &  2{,}975 &   0.948  & 339 &   0.965  \\
    2{,}000 &  1{,}078 &   0.986  &  2{,}893 &   0.982  & --- & --- \\
    \midrule
    5{,}000 & \textbf{6{,}857} & 0.775 & \textbf{11{,}178} & 0.841 & --- & --- \\
    \midrule
    10{,}000 &  6{,}273 & 0.992 &  9{,}622 & 0.989 & --- & --- \\
    13{,}000 &  5{,}278 & 0.984 &  9{,}538 & 0.979 & --- & --- \\
    50{,}000 & --- & --- &  9{,}401 & 0.984 & --- & --- \\
    \bottomrule
  \end{tabular}
\end{table}

Three features of this are worth stating explicitly.

\paragraph{The collapse is abrupt.} Dead features are roughly flat and slowly
\emph{declining} through step 2{,}000 in both long runs, then jump by a factor of
$6.4\times$ (Llama) and $3.9\times$ (Qwen) by step 5{,}000. This is not gradual
attrition. The mechanism is a ratchet: under a hard top-$k$ constraint, a feature that
falls outside the top $k$ on every batch in the window receives an encoder gradient of
exactly zero, so it cannot be pushed back into contention by training. Once the
learning rate has annealed enough that the surviving features stop reshuffling, the
set of live features is frozen.

\paragraph{Reconstruction quality does not warn you.} FVE at step 5{,}000 dips (0.775,
0.841) but this is batch noise, not signal: by step 10{,}000 both runs are at 0.99,
higher than before the collapse, and they stay there. A practitioner monitoring FVE
would see a healthy, monotonically improving run. The dictionary lost 40--60\% of its
capacity and the loss curve did not notice, because the surviving features simply
absorbed the reconstruction work.

\paragraph{Higher expansion collapses harder.} Qwen (8.0$\times$ expansion) ends with
57.4\% of its dictionary dead against Llama's 20.2\% (5.33$\times$), measured over the
full corpus. This is the expected direction---a larger dictionary relative to $d$ has
more features competing for the same $k$ slots---but with two long runs we cannot
separate it from the confound that Qwen also trained nearly four times as long. Mistral
(4.0$\times$, 1{,}000 steps) shows 0.1\% dead, which reflects only that it stopped
before the transition.

\subsection{Dense-Feature Degeneracy}
\label{sec:dense}

Sparsity of the code does not imply that the code is evenly used. With $k = 128$ of
$D = 16{,}384$, a feature would fire on 0.78\% of tokens if activation were spread
uniformly. Table~\ref{tab:dense} shows the actual distribution.

\begin{table}[t]
  \centering
  \caption{Concentration of activation mass, measured over the full 500{,}000-token
    corpus. ``Top-200 mass'' is the share of all feature activations accounted for by
    the 200 most frequently firing features---1.2\% of the dictionary.
    Data: \texttt{data/sae-analysis/corrections.json}.}
  \label{tab:dense}
  \smallskip
  \begin{tabular}{lrrrr}
    \toprule
    Model & Uniform expectation & Features firing $>$50\% of tokens & Top-200 mass \\
    \midrule
    Llama-3.2-3B & 0.78\% & 34 & 40.3\% \\
    Qwen2.5-3B   & 0.78\% & 20 & 38.6\% \\
    Mistral-7B   & 0.78\% & 34 & 57.2\% \\
    \bottomrule
  \end{tabular}
\end{table}

In every model, 1.2\% of the dictionary carries roughly two-fifths of the sparse code,
and a few dozen features fire on the majority of all tokens. A feature firing on 97\% of
tokens---the maximum we observe in both Llama and Qwen---is not a feature in any useful
sense; it is closer to a learned bias absorbed into the dictionary.

\paragraph{Consequence for feature selection.} The standard first step in feature
labelling is to take the top-$N$ features by activation frequency and inspect their
maximally activating contexts. Table~\ref{tab:dense} says that this procedure draws
almost exclusively from the degenerate head: of the 200 most frequent Llama features,
34 fire on more than half of all tokens. In separate work we built a feature-based
safety classifier on exactly this selection and obtained ROC-AUC 0.564 with a 95\%
confidence interval spanning chance. Selecting by \emph{differential} activation
between contrastive corpora, and imposing a density ceiling, are both necessary; neither
is standard practice.

\subsection{Held-Out Reconstruction}
\label{sec:heldout}

\begin{table}[t]
  \centering
  \caption{Reconstruction over the full corpus. Only Mistral has unseen data, because
    it was trained on the first 50{,}000 tokens of a 500{,}000-token dump; its held-out
    row covers tokens 50k--500k. Dead-feature counts here are features that fire nowhere
    in the corpus, a stricter criterion than the rolling window of
    Table~\ref{tab:collapse}.}
  \label{tab:heldout}
  \smallskip
  \begin{tabular}{llrrrl}
    \toprule
    Model & Split & MSE & FVE & Dead features & \\
    \midrule
    Llama-3.2-3B & in-sample & 0.00788 & 0.9865 & 3{,}309 (20.2\%) & 53 epochs \\
    Qwen2.5-3B   & in-sample & 0.27532 & 0.9839 & 9{,}398 (57.4\%) & 205 epochs \\
    Mistral-7B   & in-sample & 0.00415 & 0.9273 &     24 (0.1\%)  & 41 epochs \\
    Mistral-7B   & \textbf{held-out} & 0.00426 & \textbf{0.9252} & 26 (0.2\%) & unseen \\
    \bottomrule
  \end{tabular}
\end{table}

The Llama and Qwen SAEs each made 53 and 205 passes over their entire activation dump.
There is no unseen data for them, and at that level of repetition an SAE can in
principle memorize the corpus. We report their FVE as the best available
characterization, not as evidence of generalization.

Mistral is the exception, by accident: trained on only the first 50{,}000 tokens, it
leaves 450{,}000 unseen. Its held-out FVE is 0.9252 against an in-sample 0.9273---a gap
of 0.2 percentage points. A 1{,}000-step SAE has many problems, but overfitting its
training corpus is not among them. This is weak evidence that the other two figures are
not purely memorization; it comes from the least-trained model in the set and should
not carry much weight. Reserving a held-out split before training is the cheapest fix
available to this work.

\subsection{Cross-Architecture Feature Matching}
\label{sec:matching}

\begin{table}[t]
  \centering
  \caption{Matchable feature inventory over the 500{,}000-token corpus. ``Dead'' features
    fire nowhere; ``low support'' features fire in fewer than 50 of 1{,}000 chunks.
    Only the last column enters the matching analysis.}
  \label{tab:inventory}
  \smallskip
  \begin{tabular}{lrrrr}
    \toprule
    Model & Dictionary & Dead & Low support & \textbf{Matchable} \\
    \midrule
    Llama-3.2-3B & 16{,}384 & 3{,}309 &  4{,}771 & \textbf{8{,}304} (50.7\%) \\
    Qwen2.5-3B   & 16{,}384 & 9{,}398 &      22 & \textbf{6{,}964} (42.5\%) \\
    Mistral-7B   & 16{,}384 &     24 & 10{,}881 & \textbf{5{,}479} (33.4\%) \\
    \bottomrule
  \end{tabular}
\end{table}

Between a third and a half of each dictionary is usable (Table~\ref{tab:inventory}),
and the two failure modes are the ones diagnosed above: Qwen is half dead from
collapse, Mistral is overwhelmingly low-support because its features never specialized.

\paragraph{The null.} Table~\ref{tab:null} reports the chunk-permutation null. The
number to attend to is the mean: two features from different models with no functional
relationship, matched reciprocally over a 1{,}000-dimensional fingerprint, correlate at
0.39--0.57 on average.

\begin{table}[t]
  \centering
  \caption{Chunk-permutation null, 20 replicates. $\tau$ is set to the null 99th
    percentile. ``Expected FP'' is the number of matches a shuffled control produces
    above $\tau$. Data: \texttt{data/sae-analysis/matching-v2/matching-report.json}.}
  \label{tab:null}
  \smallskip
  \begin{tabular}{lrrrr}
    \toprule
    Pair & Null mean & Null p99 ($= \tau$) & Null matches/replicate & Expected FP \\
    \midrule
    Llama--Qwen    & 0.574 & 0.939 & 700 & 7.0 \\
    Mistral--Qwen  & 0.393 & 0.832 & 486 & 4.9 \\
    Llama--Mistral & 0.464 & 0.894 & 557 & 5.6 \\
    \bottomrule
  \end{tabular}
\end{table}

A fixed threshold of $\tau = 0.8$---the conventional choice, and the one used in an
earlier analysis of these same checkpoints---is at or below that noise floor for two of
three pairs. The problem is structural rather than a matter of picking a better number:
a match is the best of $\sim$5{,}000--8{,}000 candidates, and the expectation of an
extreme-value statistic grows with the size of the search and shrinks with fingerprint
dimensionality. At a lower resolution we also tested (100 chunks over 50{,}000 tokens,
a common configuration) the null 99th percentile reaches 0.98--0.995, leaving the
statistic with no discriminative power whatsoever.

\paragraph{Results.} Table~\ref{tab:matches} gives the counts.

\begin{table}[t]
  \centering
  \caption{Cross-architecture matches under the corrected protocol, against the
    uncalibrated protocol applied to the same checkpoints. Enrichment is observed
    matches over expected false positives.}
  \label{tab:matches}
  \smallskip
  \begin{tabular}{lrrrr}
    \toprule
    Pair & Uncalibrated & \textbf{Corrected} & Expected FP & Enrichment \\
    \midrule
    Llama--Qwen    &  5{,}923 & \textbf{22} & 7.0 & $3.1\times$ \\
    Mistral--Qwen  &  8{,}099 & \textbf{10} & 4.9 & $2.0\times$ \\
    Llama--Mistral & 12{,}174 & \textbf{12} & 5.6 & $2.1\times$ \\
    \midrule
    \textbf{Three-way universal} & \textbf{3{,}753} & \textbf{0} & 0.15 & --- \\
    \bottomrule
  \end{tabular}
\end{table}

No feature matches across all three models. The closed-triangle count is zero against a
null expectation of 0.15 per replicate. Pairwise, 10--22 pairs survive per model pair
against 5--7 expected by chance---an enrichment of 2--3$\times$, or on the order of
5--15 genuine pairs out of 5{,}000--8{,}000 matchable features (0.1--0.3\%). At these
counts Poisson uncertainty alone spans a factor of two, and we do not think the
pairwise numbers support a quantitative claim.

The three corrections compound and none is individually sufficient. Reciprocal matching
alone reduces the Llama--Mistral count from 12{,}174 to 40 at the \emph{same} threshold,
because the original figure counted a handful of Mistral features many times over. The
closed-triangle requirement removes chained triples---in the uncalibrated run, 67\% of
nominally universal triples had no verified Llama--Mistral edge, and the 3{,}753 triples
collapsed onto 644 distinct Qwen and 341 distinct Mistral features. Null calibration
removes the remainder.

\paragraph{What this does and does not show.} It does not show that these three models
fail to share representational structure. Two of the three SAEs are undertrained, all
three have degraded dictionaries, and the fingerprint statistic---mean activation over
500-token chunks---discards within-chunk structure, so a feature responding to a rare
token is indistinguishable from one responding to a chunk's topic. The negative result
is confounded with the quality of the instruments.

What it does show is methodological and transfers beyond this dataset: cross-model
feature matching by activation-pattern similarity requires a permutation null. Without
one, the search over a large dictionary manufactures high similarities from nothing, and
a threshold that reads as stringent can sit below the noise floor.

% ── 5. What Was Not Run ──────────────────────────────────────────────────────
\section{Planned Experiments Not Conducted}
\label{sec:notrun}

An earlier version of this document was framed as a three-way comparison of SAE
training objectives across three evaluation metrics, and stated results for that
comparison in its abstract. Those results did not exist. We list here what was planned
and not run, so that no reader has to reconstruct it.

\begin{itemize}
  \item \textbf{L1, Gated and JumpReLU SAEs were never trained.} Only TopK exists. Every
    claim about how objectives compare---including the claim that reconstruction quality
    mediates probe accuracy---has no data behind it.
  \item \textbf{No linear probing study.} No probes were trained on SAE features, so
    there is no probe-accuracy result and no mediation analysis.
  \item \textbf{No human interpretability study.} No annotators were recruited, no
    features were rated, and no inter-rater agreement was computed. The ethics statement
    in the earlier draft described compensation and institutional approval for a study
    that did not take place.
  \item \textbf{No steering-fidelity evaluation of SAE features.} Separate steering
    experiments exist in this workspace, but they use contrastive mean-difference
    vectors, not SAE decoder directions, and do not bear on the objective comparison.
  \item \textbf{No sparsity-level ablation.} All runs use $k = 128$; the planned
    $L_0 \in \{16, 32, 64\}$ sweep was not run.
  \item \textbf{GPT-2 Small and Pythia-1.4B were never used.} The earlier draft named
    them as the model set. No SAE was trained on either.
\end{itemize}

The objective comparison remains, in our view, the right experiment---the fragmentation
in the literature it was meant to address is real. It is simply not this paper.

% ── 6. Discussion ────────────────────────────────────────────────────────────
\section{Discussion}
\label{sec:discussion}

\paragraph{Monitor dictionary utilization, not just reconstruction.} The collapse in
\S\ref{sec:collapse} is invisible in FVE and in the loss curve. It is visible only in a
dead-feature count, and only if that count is logged during training rather than read
off the final checkpoint. Since the auxiliary loss of \citet{gao2024scaling} is designed
precisely to prevent it, the practical recommendation is simply to use it---but the
measurement matters independently, because a run using the auxiliary loss still needs
verification that it worked.

\paragraph{Report reconstruction over a corpus, not a batch.} Batch FVE at the final
step spans 0.36 across adjacent logged steps in our shortest run and 0.02 in our
longest. A single reading is a draw from that distribution. Full-corpus evaluation costs
one forward pass and removes the ambiguity entirely.

\paragraph{Expansion ratio is a hyperparameter, not a free choice.} Holding dictionary
size constant across models with different hidden dimensions---a natural way to make a
comparison ``controlled''---silently varies expansion ratio, and the model with the
highest ratio collapsed hardest. Whether expansion ratio drives collapse or merely
correlates with it here, we cannot say with two long runs.

\paragraph{Frequency-based feature selection selects the wrong features.} This follows
directly from \S\ref{sec:dense} and is, we suspect, the most transferable practical
finding. It affects any pipeline whose first step is ``take the most active features.''

\subsection{Limitations}

\textbf{Only one objective.} Everything reported here is about TopK SAEs without an
auxiliary loss. The collapse dynamics in particular are specific to hard-sparsity
training; L1 SAEs fail differently.

\textbf{Uneven and insufficient training.} One of three runs reached its target. The
Mistral SAE trained for 1{,}000 steps on a corpus ten times smaller than the others.
Cross-model differences are confounded with training length throughout, and we have
avoided attributing any difference to architecture.

\textbf{No held-out data for two of three models.} See \S\ref{sec:heldout}.

\textbf{Quantization confound.} Mistral activations come from 4-bit weights.

\textbf{Single corpus, single seed.} All results use WikiText-103 and seed 42. There is
one run per condition, so we report no variance across seeds and cannot say whether the
collapse step is stable.

\textbf{No semantic labelling.} We never establish what any feature represents. The
matching analysis asks whether activation patterns correspond, not whether the
corresponding features mean the same thing---and with zero three-way matches, the
question does not arise.

% ── 7. Conclusion ────────────────────────────────────────────────────────────
\section{Conclusion}
\label{sec:conclusion}

We trained TopK SAEs at matched sparsity and dictionary size on three model families and
instrumented the training rather than only its endpoint. Without an auxiliary revival
loss, the dictionary collapses abruptly between steps 2{,}000 and 5{,}000, losing
40--60\% of its capacity while reconstruction quality continues to improve---a failure
mode that the usual metrics conceal. What survives is dominated by dense directions:
1.2\% of the dictionary carries roughly 40\% of the sparse code, which invalidates
frequency-based feature selection.

Cross-architecture feature matching finds nothing that survives a permutation null. The
conventional protocol applied to these same checkpoints reports 3{,}753 shared features;
the corrected count is zero. The gap is entirely attributable to three procedural
choices---many-to-one matching, chained rather than closed triples, and an uncalibrated
threshold---each of which is easy to make and none of which is visible in the output.

We report the objective comparison this paper was originally framed around as not
conducted.

% ── Reproducibility ──────────────────────────────────────────────────────────
\section*{Reproducibility Statement}

All SAE training code, activation collection scripts, training logs, checkpoints, and
analysis scripts are available at
\url{https://github.com/jmelton-research/sae-crossarch}. All runs used seed 42.
Training was conducted on Apple Silicon using the MLX framework. The corrected matching
analysis is \texttt{data/sae-analysis/cross\_arch\_matching\_v2.py} and the superseded
version is retained as \texttt{cross\_arch\_matching.py} for comparison; reconstruction
and null statistics are regenerated by
\texttt{data/sae-analysis/recompute\_corrections.py}.

\section*{Ethics Statement}

This work involves no human subjects and no personal data; all activations are derived
from WikiText-103, a public corpus of Wikipedia text. An earlier draft of this document
contained an ethics statement describing annotator compensation and institutional
approval for a human interpretability study that was never conducted; that statement has
been removed.

SAEs and the feature-matching methods studied here have applications in interpretability
and safety research and potential for misuse in targeted manipulation of model behaviour.
Our results are negative and release no feature directions of either kind.

\bibliography{references}
\bibliographystyle{plainnat}

\end{document}
